\subsection{Frozen lower forests and a shared-query retry}
\label{marked-cut:lower-forests}

A complete top reduction need not leave a cancellable presentation, and
neither need a particular zero/one forest. The experimental lower-forest
probe in the preceding section isolated six source cores for which a
maximum-passage forest permits complete raw cancellation. Four already
extended to whole-cut certificates. Two failed at the bundle kernel because
that consumer silently reconstructed the historical forest. This section
makes the choice explicit throughout the finite computation and extends
independent kernel replay to authenticate it. The change broadens a
sufficient terminal criterion; it supplies no negative topological verdict.

\paragraph{The lower graph and its geometric meaning.}
Let $G$ have one vertex for each source zero-handle and one edge for each
source one-handle, with the two attaching ends as endpoints. Loops and
parallel edges retain their source identifiers. A lower forest $F$ is a
spanning tree in every connected component of $G$, including an empty tree
at an isolated vertex. Contracting its handles successively amalgamates
zero-handles along one-handles joining distinct balls. Each amalgamation
is a ball. The remaining one-handles therefore give the usual handlebody
spine on one zero-handle per component. No loop may belong to $F$.

The original face-sector charts still locate every attachment. To count
intersections with a remaining one-handle cocore, trace the source attaching
cycle and omit precisely the flags on $F$. This operation uses literal,
unsigned passages. In particular, oppositely oriented passages are counted
twice, not cancelled. Changing $F$ changes this geometric disk system and
can change which raw singleton cancellations exist, although the underlying
cut manifold has not changed. The argument certifies the local contraction
and the ensuing supported operations; it is not a general algorithm for
transporting arbitrary boundary patterns through a hierarchy.

\begin{proposition}[Maximum original passage weight]
For each source one-handle $e$, let $w_e$ count its occurrences over all
original source two-handle attaching cycles, before any top cancellation or
bundle absorption. A maximum-weight spanning forest $F$ minimizes
\[
 L(F)=\sum_{e\notin F} w_e
\]
among all lower forests of $G$. The producer constructs such a forest by
processing edges in decreasing $(w_e,-e)$ order and accepting exactly those
joining distinct current components.
\end{proposition}
\begin{proof}
The sum of all weights is fixed, so minimizing $L$ is equivalent to
maximizing $\sum_{e\in F}w_e$. Suppose an optimal spanning forest contains
the already accepted edges but omits the next accepted edge $e$. Adding
$e$ creates a cycle in its source connected component. Some edge $f$ on
that cycle crosses between distinct current components. Such an edge
cannot have been rejected earlier: a rejection would mean its endpoints
were already connected, and connections persist. Thus $w_f\leq w_e$.
Exchange $f$ for $e$ without decreasing weight or removing an earlier
accepted edge. Induction proves optimality. Scanning all edges gives a
spanning forest because any unconnected source edge would have been
accepted. The argument applies componentwise, ignores loops correctly,
and allows parallel edges.
\end{proof}

The weights are exact counts, not the saturated $0,1,2$ values used by the
cancellation planner. These are different tasks: truncation preserves pivot
admissibility but need not preserve an additive optimization objective.
For a core with $V$ zero-handles, $E$ one-handles and $P$ original attaching
flags, weight construction takes $O(P)$ increments, sorting takes
$O(E\log(E+1))$ comparisons, and union by size with path compression takes
$O(E\alpha(V+E))$ amortized disjoint-set operations. Storage is $O(V+E)$
beyond the source core. With source size $O(T)$, this is a polynomial
finite stage independent of normal sheet multiplicity. Arithmetic costs
include the bits of the counts and identifiers; the graph-stage bound is
not a bound for orbit discovery or for a full recognition hierarchy.

\paragraph{One choice, retained by every consumer.}
The implementation selects the alternative forest on the \emph{original}
source core. A relative two/three reduction leaves all zero/one handles
and their incidences unchanged. Disk-bundle absorption adds two-handles
only. Consequently the same lower forest is valid for both resulting
cores. The following consumers receive it without recomputation: the
cost calculation for the top forest, the source raw-cancellation precheck,
the augmented core certificate, core-side annulus transport through that
certificate, and the whole bundle-boundary kernel. Port owners use the
least source zero-handle identifier of their connected component, so
changing contraction order does not relabel connected components.

Freezing the choice is essential to the implementation contract. Computing
one forest for the precheck and another for the kernel tests different raw
matrices; the former's success cannot justify the latter's cancellation.
It also prevents a purported fixed-cost optimization from changing its
objective when top rows are removed. This two-stage choice is not a joint
optimization over primal and dual trees: the lower objective includes all
original rows, while the top objective subsequently minimizes retained
passages relative to that fixed lower forest. Neither objective proves
existence of a complete raw cancellation schedule across all forests.

\paragraph{An independently authenticated kernel certificate.}
The existing core checker already accepts arbitrary spanning lower forests.
The bundle checker previously reconstructed one canonical forest. The new
\path{bundle-boundary-kernel-v3} format has the same fields but interprets
its forest as supplied evidence. Independent replay decodes exact integer
identifiers, rejects repeated edges, out-of-range identifiers and cycles,
and checks that every source edge has endpoints in the same resulting
component. The last condition proves componentwise spanning. Component
labels are reconstructed by explicit label merging, without importing the
producer's selector or its disjoint-set helper.

Replay then reconstructs source occurrence rows, actual annular crossing
curves, dedicated bundle columns, connecting-handle forest and census.
The existing independent column timestamp checker verifies the cancellation
schedule. It never trusts a producer weight, an optimization claim or a
reported census. Legacy v1 and v2 certificates retain their canonical
forest semantics, including the dense v1 replay path. The outer whole-cut
formats v1--v6 require no change. Only an explicitly supplied lower forest
requests a v3 kernel proof; the historical source path emits the same
certificate as before.

\begin{proposition}[Source-first preservation and shared discovery]
The adaptive policy preserves every uninterrupted positive source-policy
result and its certificate. It performs at most two finite forest attempts
and at most one set of source and boundary orbit queries. Any successful
alternative result is justified by independent source-relative replay.
\end{proposition}
\begin{proof}
The first attempt uses the historical forest and returns immediately on
success. On a method refusal the second attempt uses maximum original
passage weight. If the selected edge set equals the historical set, it is
skipped: the raw matrix, component labels and fixed top costs are identical.
The source parallelity packet is computed before either attempt. Boundary
discovery is lazy and caches its result, so a second attempt either reuses
the first packet or causes its first computation. A boundary-budget refusal
returns immediately; changing the finite forest cannot repair that same
query allowance. Source and boundary cycles therefore share one budget,
with no retry reset. Both attempts propagate cancellation callbacks.
Soundness of a positive alternative follows from forest validation and the
same local geometric operations verified for a source-policy proof.
\end{proof}

The API exposes \path{lower_forest='adaptive'} by default, with explicit
\path{'source'}, \path{'max_passages'} and \path{'reverse_id'} policies for
reproduction and experiments. Reverse identifier order is not an additional
default attempt. The default has four source queries and the previously
chosen one-query moment or two-query mask boundary backend, regardless of
whether the source precheck or a later kernel step causes the retry.
It can spend more finite time on an unresolved input; no elapsed-time
improvement follows merely from having a second choice. The precheck still
means only that this particular raw presentation cannot be completely
cancelled by the supported rules. A refusal never proves knottedness.

Useful-surface discovery, general boundary-pattern transport, global
hierarchy and restart accounting, and a native bit bound for orbit traces
remain separate unresolved obligations. In particular this bounded retry,
its spanning-tree optimality, and any extra finite-corpus successes do not
establish quasi-polynomial unknot recognition.
